%%%%%%%%%%%%%%%%%%%%%%%%%%%%%%%%%%%%%%%%%%%%%%%%%%%%%%%%%%%%%%%%%%%%%%%%%%%%%%%
%
% witseiepaper-2005.tex
%
%                       Ken Nixon (12 October 2005)
%
%                       Sample Paper for ELEN417/455 2005
%
%%%%%%%%%%%%%%%%%%%%%%%%%%%%%%%%%%%%%%%%%%%%%%%%%%%%%%%%%%%%%%%%%%%%%%%%%%%%%%%%

\documentclass[10pt,twocolumn]{witseiepaper}
%
% All KJN's macros and goodies (some shameless borrowing from SPL)
\usepackage{KJN}
\usepackage[super]{nth}
\usepackage{subcaption}
\usepackage{listings}
\usepackage{amsmath}
\usepackage{epstopdf}
\usepackage{xcolor}
\usepackage{textcomp}
\usepackage{listings}
\usepackage{alltt}
%\usepackage{matlab-prettifier}
\usepackage{graphicx}
\usepackage{changes}
\usepackage{makecell}
\usepackage{verbatim}
\usepackage{algorithm,algpseudocode}
\usepackage{balance}
\usepackage{pdfpages}
\usepackage{ragged2e}
\usepackage{color} %red, green, blue, yellow, cyan, magenta, black, white
\definecolor{mygreen}{RGB}{28,172,0} % color values Red, Green, Blue
\definecolor{mylilas}{RGB}{170,55,241}
%\usepackage{flafter}

%
% PDF Info
%
\ifpdf
\pdfinfo{
/Title (INSTRUCTIONS AND STYLE GUIDELINES FOR THE PREPARATION OF FINAL YEAR LABORATORY PROJECT PAPERS : 2005 VERSION)
/Author (Ken J Nixon)
/CreationDate (D:200309251200)
/ModDate (D:200510121530)
/Subject (ELEN417/455 Paper Format, 2005)
/Keywords (ELEN417, ELEN455, paper, instructions, style guidelines, laboratory project)
}
\fi

%%%%%%%%%%%%%%%%%%%%%%%%%%%%%%%%%%%%%%%%%%%%%%%%%%%%%%%%%%%%%%%%%%%%%%%%%%%%%%%
\begin{document}
	
\begin{titlepage}
	
	\newcommand{\HRule}{\rule{\linewidth}{0.3mm}} % Defines a new command for the horizontal lines, change thickness here
	
	\center % Center everything on the page
	
	%----------------------------------------------------------------------------------------
	%	HEADING SECTIONS
	%----------------------------------------------------------------------------------------
		\includegraphics[width=0.3\textwidth]{EIE.png}\\[1cm] % Include a department/university logo - this will require the graphicx package
	
	%----------------------------------------------------------------------------------------
	\textsc{\LARGE University of the Witwatersrand } \\[0.1cm] % Name of your university/college
	\textsc{\LARGE School of Electrical and Information Engineering }\\[1cm] % Major heading such as course name
	\textsc{\Large ELEN4020: Data Intensive Computing}\\[1.5cm] % Minor heading such as course title
	
	%----------------------------------------------------------------------------------------
	%	TITLE SECTION
	%----------------------------------------------------------------------------------------
	
	\HRule \\[0.4cm]
	{ \huge \bfseries Laboratory Exercise 1} \\[0.4cm] % Title of your document
		\HRule \\[1.5cm]

%----------------------------------------------------------------------------------------
%	AUTHOR SECTION
%----------------------------------------------------------------------------------------
\textsc{\Large 	\emph{Authors:} } \\[0.1cm]	 


\begin{minipage}{0.4\textwidth}
	\begin{flushleft} \large
		%			\emph{Author:} \\
		Kayla-Jade Butkow \\ 714227 % Your name
	\end{flushleft}
\end{minipage}
~
\begin{minipage}{0.4\textwidth}
	\begin{flushright} \large
		%	\emph{Author:}\\
		Jared Ping \\ 704447
	\end{flushright}
\end{minipage}\\[1cm]

\begin{minipage}{0.4\textwidth}
	\begin{flushleft} \large
		%		\emph{Author:}\\
		Lara Timm \\ 704157
	\end{flushleft}
\end{minipage}
~
\begin{minipage}{0.4\textwidth}
	\begin{flushright} \large
		%		\emph{Author:} \\
		Matthew van Rooyen \\ 706692
	\end{flushright}
\end{minipage}\\[1cm]
		
		
	
{\large Date Handed In: \nth{23} February, 2018}\\[1cm] 
	
\end{titlepage}


%%%%%%%%%%%%%%%%%%%%%%%%%%%%%%%%%%%%%%%%%%%%%%%%%%%%%%%%%%%%%%%%%%%%%%%%%%%%%%%
\thispagestyle{empty}
\pagestyle{plain}
\setcounter{page}{1}

\section{Introduction}


\section{PThread and OpenMP Libraries}

\subsection{PThread Library}
PThreads is a standardized model for dividing a program into parallel tasks \cite{pthreads}. PThreads was defined by the IEEE POSIX operating system interface standards \cite{pthreads}. The PThreads library specifies the interface to manage the actions required by threads \cite{pthreadVSopen}.

\subsubsection{Disadvantages}
In order to use the PThreads library, the code must be written specifically for the library \cite{pthreadVSopen}. This involves the use of PThread specific functions and data structures. The implication of this is that once the library has been used, the application is inherently casted as a threaded model \cite{pthreadVSopen}. Furthermore, since the API is very low level, large amounts of complex code is required to perform simple threading tasks \cite{pthreadVSopen}. The developer is also required to control all stages of the threading process, such as implicitly stating the number of threads and terminating them after the parallel block has executed \cite{pthreadVSopen}.

\subsubsection{Advantages}
However, the use of the PThreads library allows the developer extensive control over all aspects of the program.

\section{Conclusion}

\bibliographystyle{witseie}
\bibliography{dataLab1}

\end{document}

" vim: ts=4
" vim: tw=78
" vim: autoindent
" vim: shiftwidth=4
